\documentclass[10pt,twocolumn,a4paper]{article}

% Language setting
\usepackage[vietnamese]{babel}
\usepackage[utf8]{inputenc}

% Set page size and margins - Tighter for 2-column layout
\usepackage[top=2cm, bottom=2cm, left=1.5cm, right=1.5cm]{geometry}

% Font setting - Times New Roman style
\usepackage{mathptmx}

% Useful packages
\usepackage{amsmath}
\usepackage{graphicx}
\usepackage{booktabs}
\usepackage[colorlinks=true, allcolors=blue]{hyperref}
\usepackage{natbib}

\title{\Large \textbf{Các Phương pháp Tối ưu hóa Cache}}
\author{Nhóm nghiên cứu: [Tên thành viên 1, 2, 3, 4]}
\date{\today}

\begin{document}

\maketitle

\begin{abstract}
Cache đóng vai trò then chốt trong việc thu hẹp khoảng cách hiệu suất giữa bộ xử lý và bộ nhớ chính. Trong kỷ nguyên dữ liệu lớn và tính toán hiệu năng cao, việc tối ưu hóa cache không chỉ dừng lại ở các tầng phần cứng gần CPU mà còn mở rộng ra các kỹ thuật quản lý bộ nhớ ảo và dự đoán dữ liệu. Bài tiểu luận này trình bày kết quả nghiên cứu của nhóm về bốn phương pháp tối ưu hóa cache quan trọng: (1) Sử dụng bộ nhớ ảo để mở rộng dung lượng RAM, (2) Kiến trúc cache đa cấp, (3) Kỹ thuật nạp trước (Precaching), và (4) [Chủ đề phần 4]. Mục tiêu của nghiên cứu là phân tích cơ chế hoạt động, ưu điểm và thách thức của từng phương pháp trong việc cải thiện tổng thể hiệu năng hệ thống.
\end{abstract}

\section{Giới thiệu}

Trong kiến trúc máy tính hiện đại, bộ nhớ cache là một thành phần không thể thiếu nhằm giải quyết vấn đề "bức tường bộ nhớ" (memory wall) - hiện tượng tốc độ của bộ xử lý (CPU) tăng trưởng nhanh hơn nhiều so với tốc độ truy cập bộ nhớ chính (DRAM). Một hệ thống cache hiệu quả có thể giảm thiểu độ trễ truy cập, tăng băng thông dữ liệu và tối ưu hóa mức tiêu thụ năng lượng.

Nhóm nghiên cứu của chúng tôi tập trung vào bốn khía cạnh tối ưu hóa khác nhau, từ mức độ hệ điều hành đến mức độ phần cứng vi kiến trúc. Mỗi thành viên chịu trách nhiệm phân tích sâu một kỹ thuật cụ thể, đảm bảo cái nhìn toàn diện về cách thức dữ liệu được lưu trữ và truy xuất hiệu quả nhất.

\section{Phần 1: Bộ nhớ ảo và Mở rộng RAM bằng Ổ cứng}

\subsection{Khái niệm Bộ nhớ ảo (Virtual Memory)}
Bộ nhớ ảo là một kỹ thuật quản lý bộ nhớ được thực hiện bởi sự phối hợp giữa phần cứng (MMU - Memory Management Unit) và phần mềm (Hệ điều hành). Nó cho phép hệ thống sử dụng một không gian địa chỉ logic lớn hơn nhiều so với dung lượng vật lý thực tế của RAM. Về bản chất, bộ nhớ ảo hoạt động như một tầng cache cho ổ cứng, nơi các trang (pages) dữ liệu được hoán đổi linh hoạt.

\subsection{Hạn chế của cơ chế Swap truyền thống}
Cơ chế hoán đổi (Swap) truyền thống hoạt động ở mức độ \textbf{trang (page granularity)} - thường là 4KB. Điều này dẫn đến hai vấn đề lớn khi sử dụng với ổ cứng thể rắn (SSD):
\begin{itemize}
    \item \textbf{Write Amplification}: Khi chỉ một đối tượng nhỏ trong trang bị thay đổi, toàn bộ trang 4KB vẫn phải ghi lại xuống SSD, làm giảm hiệu năng và tuổi thọ chip Flash.
    \item \textbf{Lãng phí RAM}: Các trang nạp vào RAM thường chứa nhiều đối tượng "lạnh" không cần thiết.
\end{itemize}

\subsection{Giải pháp Hybrid SSD/RAM: Hệ thống SSDAlloc}
Nghiên cứu hiện đại (điển hình như hệ thống \textbf{SSDAlloc} \cite{badam2011ssdalloc} từ Đại học Princeton) đề xuất việc coi SSD như một phần mở rộng trực tiếp của RAM thay vì chỉ là thiết bị lưu trữ thứ cấp. Các đặc điểm chính bao gồm:
\begin{itemize}
    \item \textbf{Object Granularity}: Thay vì quản lý theo trang 4KB, hệ thống quản lý dữ liệu ở mức độ đối tượng ứng dụng.
    \item \textbf{Mô hình Object-Per-Page (OPP)}: Mỗi đối tượng nhỏ được gán vào một trang bộ nhớ ảo riêng biệt.
    \item \textbf{RAM Object Cache}: Sử dụng phần lớn RAM làm bộ nhớ đệm cho các đối tượng "nóng".
\end{itemize}

\subsection{Tối ưu hóa Hiệu năng và Tuổi thọ}
Việc kết hợp SSD và RAM theo mô hình hybrid mang lại những cải thiện vượt bậc:
\begin{itemize}
    \item \textbf{Hiệu năng}: Tăng băng thông xử lý lên đến 17 lần so với cơ chế Swap thông thường của Linux.
    \item \textbf{Độ bền}: Tuổi thọ của SSD được kéo dài đáng kể thông qua việc giảm thiểu số lần ghi không cần thiết.
\end{itemize}

\subsection{Ứng dụng trong Hệ thống Phân tán}
Trong các hệ thống phân tán, bộ nhớ ảo hybrid cho phép các nút xử lý các tập dữ liệu cực lớn ngay trên một máy chủ đơn lẻ mà không cần đầu tư quá nhiều vào RAM vật lý.

\section{Phần 2: Cache Đa cấp (Multilevel Cache)}

\subsection{Cấu trúc Phân cấp Bộ nhớ}
Để tối ưu hóa cả tốc độ và dung lượng, các kiến trúc sư máy tính sử dụng nhiều tầng cache giữa CPU và RAM:
\begin{itemize}
    \item \textbf{L1 Cache}: Nhỏ nhất, nhanh nhất, nằm ngay trong lõi CPU.
    \item \textbf{L2 Cache}: Dung lượng lớn hơn, độ trễ cao hơn L1 một chút.
    \item \textbf{L3 Cache}: Thường được chia sẻ giữa tất cả các lõi.
\end{itemize}

\subsection{Chính sách Inclusivity (Bao hàm)}
Việc quản lý dữ liệu giữa các cấp cache tuân theo các chính sách chính: \textit{Inclusive}, \textit{Exclusive}, và \textit{Non-inclusive}.

\subsection{Tối ưu hóa Hiệu năng qua Cache Đa cấp}
Hiệu năng được đo bằng AMAT:
\[
\text{AMAT} = \text{T}_1 + \text{MR}_1 \times (\text{T}_2 + \text{MR}_2 \times \text{T}_{\text{memory}})
\]

\subsection{Xu hướng hiện đại: L4 Cache và HBM}
Sử dụng bộ nhớ băng thông cao (HBM) ngay trên đế chip để tạo ra một tầng đệm khổng lồ.

\section{Phần 3: Nạp trước dữ liệu (Precaching)}

\subsection{Nguyên lý của Precaching}
Dự đoán các dữ liệu hoặc lệnh mà CPU sẽ cần trong tương lai gần và nạp chúng vào cache trước khi chúng thực sự được yêu cầu.

\subsection{Phân loại Kỹ thuật Prefetching}
\begin{enumerate}
    \item \textbf{Hardware Prefetching}: \textit{Stride Prefetcher}, \textit{Stream Buffer}.
    \item \textbf{Software Prefetching}: Chèn các lệnh gợi ý vào mã nguồn. 
\end{enumerate}

\subsection{Lợi ích và Rủi ro}
\textbf{Lợi ích}: Giảm thời gian chờ của CPU.
\textbf{Rủi ro}: \textit{Cache Pollution} (làm bẩn cache) nếu dự đoán sai.

\subsection{Ứng dụng trong Web Caching và CDN}
Các trình duyệt và CDN nạp trước tài nguyên để tăng tốc độ truy cập cho người dùng.

\section{Phần 4: [Tên chủ đề phần 4]}

\textit{(Nội dung phần này hiện đang để trống)}

\section{Kết luận}

Tối ưu hóa cache là một quá trình đa tầng. Việc kết hợp linh hoạt các phương pháp giúp hệ thống hiện đại đạt được hiệu năng tối ưu, đáp ứng nhu cầu của các ứng dụng phức tạp.

\bibliography{sample}
\bibliographystyle{plain}

\end{document}
